\documentclass{article}
\usepackage{graphicx}
\usepackage{amssymb}
\usepackage{amsmath}
\usepackage{commath}
\usepackage{amsthm}
\newtheorem*{remark}{Remark}
\usepackage{enumitem}
\usepackage{outlines}
\setlist[enumerate,1]{label=\roman*)}
\setlist[enumerate,2]{label=\alph*)}

%PREAMBLE
\title{Report of PhD activities in the academic year 2017 / 2018}
\author{Nora Juhasz}

%CONTENT
\begin{document}

\maketitle
\newpage

\begin{abstract}

This document describes the achievements, current state and future plans of the research project, moreover it contains the courses and conferences attended in the 2017 / 2018 academic year.\\

\end{abstract}
\newpage

\tableofcontents
\newpage

%RESULTS
\section{Current state of the research topic - Atmosphere modelling}

\subsection{Achievements}

\subsection{Outlook}

%SEMINARS, SUMMER SCHOOLS
\section{Courses, summer schools and seminars}

\subsection{Courses}

\subsubsection{Atmosphere dynamics - Prof G. Curci, Prof. G. Redaelli, University of L'Aquila}

Physics of the atmosphere.

\subsubsection{Numerical methods for hyperbolic problems - Prof. G. Russo (University of Catania), GSSI}

This course gave an overview on shock capturing schemes for the numerical solution of systems of conservation and balance laws. The emphasis was on the construction of high order finite volume and finite difference schemes, which are based on the three main ingredients: numerical flux, high order non oscillatory reconstruction, and high order integration in time. The first part of the course illustrated some models in which the nonlinearity produces shocks that propagate in time. Simple three point schemes, such as upwind, Lax-Friedrichs, and Lax-Wendroff was reviewed, and their behaviour on linear equations and systems analysed, before the illustration of more complex methods for the quasilinear systems. The step by step construction of the schemes was illustrated in lab sessions, and the methods were applied to concrete examples of physical interest. The last part of the course was devoted to systems with source term.

\subsection{Conferences, seminars, summer schools}

\subsubsection{Intensive Program on Fluids and Waves, GSSI}

This training activity within the European Innovating Training Network ModCompShock surveyed recent results in the analysis of fluids and waves with emphasis on modelling and computational methods as well as its theoretical aspects.

\begin{outline}[enumerate]
  \1 Week 1. Workshop/School in Numerics for Fluids and Waves
  \1 Week 2. Mini courses by J. Bedrossian. and D. Cordoba
  \1 Week 3. Mini courses by C. Bardos, Z. Grujic, J. Malek
  \1 Week 4. Workshop on Mathematical fine structures in fluid dynamics
\end{outline}

\subsubsection{Women in applied and computational Mathematics, GSSI}

This three-day event was designed for researchers challenged by the application of various techniques in modern analysis and computational mathematics to intriguing problems arising in physics, material science and biology. The topic and the results of this present research project were represented - together with the work of many young female researchers - in the poster session.

\subsubsection{Evans function methods in the stability analysis of periodic wavetrains, Prof. R. Plaza}

This lecture series provided an introduction to recent techniques in the spectral stability analysis of spatially periodic travelling waves which involve the periodic Evans function introduced by Gardner.

\subsubsection{On the spectral stability of traveling fronts for reaction diffusion-degenerate equations, Prof. R. Plaza}

\subsubsection{Secondment at the University of Sussex, Prof. O. Lakkis}

The aim of this 6-week secondment was to create both the theoretical and technical foundations for numerical simulations in order to visualise the results of the project. Results, activities:

\begin{outline}[enumerate]

\1 Software packages such as Docker, Anaconda, ParaView, FEniCS and Alberta have been installed
\1 Writing scripts to login / connect to the HPC server at the Sussex University
\1 Getting to know FEniCS: tutorial files and the FEniCS version of the "hello world" programs
\1 Specific theoretical details that arise when implementing this project - the loss of incompressibility, Lagrange multipliers.

\end{outline}

\subsubsection{Prague}

Prague desc.

\nocite{*}
\bibliographystyle{apalike}
\bibliography{merged_U_C_eqs}

\end{document}
